% ============================================================
%  Annals of Nuclear Energy / Elsevier elsarticle submission source.
%  Generated from manuscript/paper_final_en.tex without changing scientific data.
% ============================================================
\documentclass[preprint,12pt]{elsarticle}

\usepackage{graphicx,booktabs,amsmath,amssymb,enumitem,multirow,float,array}
\usepackage[hidelinks]{hyperref}
\graphicspath{{figures/}{../figures/publication/}{../figures/hardsphere_pf_sweep/}{../figures/hardsphere_pf618_seed_replicates/}{../figures/hardsphere_pf618_spectrum/}{../figures/hardsphere_spatial_diagnostics/}{../figures/mechanism_pf62/}}
\setlength{\emergencystretch}{2em}
\biboptions{numbers,sort&compress}

\journal{Annals of Nuclear Energy}

\begin{document}

\begin{frontmatter}

\title{Pebble-scale topology effects on homogenised neutronics of Li$_4$SiO$_4$/Be breeder beds}

\author[a,b]{Jian Wang\corref{cor1}}
\ead{wjfttt@mail.ustc.edu.cn}
\author[a]{Hang Zhang}
\author[a]{Haoxi Wang}
\author[a]{Zhiyuan Liu}
\author[a]{Haishun Deng}
\author[a,b]{Wei Wen}
\author[a]{Gang Shen}

\cortext[cor1]{Corresponding author.}
\address[a]{Anhui University of Science and Technology, Huainan, Anhui, China}
\address[b]{Institute of Plasma Physics, Hefei Institutes of Physical Science, Chinese Academy of Sciences, Hefei, Anhui, China}

\begin{abstract}
Volume-homogenised breeder-bed models remove the connected He void topology
that controls neutron free paths, first-collision ordering and fast-neutron
exposure of Be multiplier pebbles.  This study compares material-equivalent
homogenised and geometry-resolved Li$_4$SiO$_4$/Be representative volume
elements (RVEs) in OpenMC using strict non-overlapping hard-sphere packings.
For the audited baseline at packing fraction (PF) = 0.618, with reflective
lateral boundaries and axial vacuum boundaries, the resolved geometry increases
the local $^6$Li(n,t) tritium-production response (TPR$_6$) by
$+17.47\pm0.71$ \%, total H3-production by $+29.13\pm0.25$ \% and
$^9$Be(n,2n) by $+45.80\pm0.11$ \%, while reducing the OpenMC neutron-mode
heating score by $-18.49\pm0.04$ \% and increasing track-length flux by
$+23.08\pm0.09$ \%.  A PF sweep from 0.56 to 0.618 shows decreasing
topology-induced shifts with increasing PF, and three independent PF=0.618
realisations preserve the sign and order of the bias.  A 50-group spectrum
diagnostic gives a resolved/homogenised flux ratio of 1.221 above the
approximate $^9$Be(n,2n) threshold, supporting void-channel streaming and
delayed solid collisions.  ENDF/B-VIII.0 and FENDL-3.2c reruns preserve the
same bias trend, while an all-reflective check shows that neutron return
changes the partitioning of the local response.  A photon-transport rerun gives
a heating shift of $-17.70\pm0.17$ \%, close to the neutron-mode result.  These
are source-normalised local RVE responses, not full-blanket TBR margins.
\end{abstract}

\begin{keyword}
Fusion breeder blanket \sep OpenMC \sep Pebble-bed microstructure \sep $^6$Li(n,t) tritium-production response \sep Beryllium multiplier \sep Homogenisation bias
\end{keyword}

\end{frontmatter}

\section{Introduction}
% ============================================================

Solid ceramic breeder blankets, in which millimetre-scale Li$_4$SiO$_4$ or
Li$_2$TiO$_3$ pebbles are mixed with beryllium multiplier pebbles and cooled by
helium purge gas, are central to several International Thermonuclear
Experimental Reactor (ITER) test blanket module (TBM), China Fusion Engineering
Test Reactor (CFETR) and demonstration fusion power plant (DEMO) blanket
concepts~\cite{giancarli2018,fischer2015,wang2020}.  Neutronic analyses of
these systems commonly replace the granular breeder bed with a
volume-homogenised mixture~\cite{fischer2015,shen2020,wang2020}.  This
approximation is attractive for full-device Monte Carlo transport because it
reduces geometric complexity and makes large blanket sectors computationally
tractable.  However, homogenisation also removes the inter-pebble He pore
network, local packing-fraction variations and discrete breeder--multiplier
interfaces.  These features influence the early collision history of
14.07\,MeV deuterium--tritium (D--T) source neutrons.  The loss of
pebble-scale topology is therefore not merely a modelling convenience issue,
but a neutronic accuracy issue for local $^6$Li(n,t) tritium-production response (TPR$_6$),
$^9$Be(n,2n) multiplication, neutron flux and heating response.

The physical basis for this concern is well established.  Random sphere
packings exhibit wall-induced packing-fraction oscillations and approach their
interior packing state only after several particle diameters
~\cite{bernal64,jodrey85,gong2017,zhang2022}.  In a solid breeder zone, this
heterogeneity changes neutron free paths, first-collision locations and the
competition among threshold multiplication in Be, absorption in $^6$Li-bearing
breeder material and energy deposition in the solid phases.  When the same
material inventory is homogenised before solving transport, these energy- and
space-dependent processes are collapsed into a single macroscopic cross
section.  The key issue is therefore to determine when a homogenised
description preserves the relevant reaction balance and when explicit pebble
geometry is required.

Previous pebble-bed neutronics studies have shown that discrete geometry can
affect tritium production and local reaction distributions.  For example, Zhang
and Guo~\cite{zhang2025} compared homogenised, regular-packing and
random-packing models for Li$_2$O mock-up and helium-cooled solid-breeder
blanket calculations.  Their random-packing model reduced the average C/E
deviation of diagnostic-pellet tritium-production rates by approximately
20 \%, and their Li$_4$SiO$_4$ blanket analysis showed that a small
zone-integrated deviation can coexist with appreciable local deviations near
pebble-bed boundaries.  These studies establish the importance of explicit
particle geometry, but four gaps remain for
mechanism-oriented RVE analysis.  First, the physical interpretation of the
bias is often limited to integral tritium production, with less attention to
neutron spectrum evolution, $^9$Be(n,2n) threshold reactions, local heating and
target-normalised reaction rates.  Second, many comparisons use one packing or
one material composition, so the sensitivity to PF, Be loading and independent
packing realisations is not fully quantified.  Third, the distinction between a
local $^6$Li(n,t) response, total H3-production response and full-blanket TBR is
often not stated sharply enough for reviewer-proof interpretation.  Fourth, the
magnitude of the bias depends strongly on boundary conditions.  Results
obtained in leakage-dominated local models therefore cannot be transferred
directly to reflected, multilayer engineering blankets without qualification.

The underlying neutronics ingredients of the present calculation are supported
by external fusion-blanket benchmark studies.  The Frascati Neutron Generator
(FNG) HCPB test blanket module mock-up measured tritium production rates and
neutron flux in a Be/Li-bearing blanket environment under 14 MeV D--T neutron
irradiation, providing a reference experiment for nuclear data and transport
methods in ceramic breeder blanket analyses~\cite{batistoni2007}.  More recent
water-cooled ceramic breeder (WCCB) mock-up experiments further demonstrated
tritium-production measurements in CFETR-relevant ceramic breeder blanket
configurations under D--T neutron irradiation~\cite{zhu2021}.  In parallel,
OpenMC has been compared with MCNP, Serpent and experimental fusion benchmarks,
including FNG-HCPB-type responses, with good agreement for tritium-production
and activation responses in the benchmark problems considered~\cite{valentine2022}.
As an executable check of the transport environment used here, we also
reproduced the OKTAVIAN aluminium-sphere neutron-leakage benchmark from the
SINBAD fusion benchmark suite~\cite{kodeli2021,konno2010}.  The present OpenMC
calculation gave a median calculation-to-experiment (C/E) ratio of 1.213, a
geometric-mean C/E of 1.046 and 85.7 \% of valid energy groups within a factor
of two of the measured leakage spectrum.  The corresponding public reference
calculations give 87.3 \% for OpenMC 0.14.0 with ENDF/B-VIII.0 and 88.8 \% for
MCNP 6.1.1b with ENDF/B-VIII.0.  This benchmark does not validate the
Li$_4$SiO$_4$/Be RVE geometry directly, but it verifies that the local OpenMC
installation, D--T source treatment and source-normalised leakage-spectrum
tallying reproduce a standard fusion-neutronics experiment at the same overall
accuracy level as published code-benchmark calculations.
These studies support the D--T source normalisation, Li-bearing
tritium-production tallies, Be multiplier physics and continuous-energy Monte
Carlo workflow used here, but they do not constitute a direct validation of the
Li$_4$SiO$_4$/Be RVE topology.  The present work therefore isolates a different
question: the local homogenisation bias
introduced by replacing a resolved Li$_4$SiO$_4$/Be/He pebble topology with a
material-equivalent uniform medium.

The central scientific question of this work is therefore how large a local
homogenisation bias is introduced when a Li$_4$SiO$_4$/Be/He pebble-bed topology
is replaced by a material-equivalent uniform medium, which transport mechanism
controls the bias, and how robust this bias is to controlled changes in PF,
multiplier inventory and hard-sphere realisation.  To address this question, we
use a paired OpenMC workflow in which each representative volume element (RVE)
is represented both as a homogenised material and as explicit constructive solid
geometry (CSG) spheres.  The two descriptions conserve the same breeder,
multiplier and purge-gas inventory.  Differences in TPR$_6$, total
H3-production response, $^9$Be(n,2n), flux and heating can therefore be
attributed to geometric topology, neutron free path and collision ordering
rather than to changes in material composition.

This study makes three contributions.  First, it establishes a strict
hard-sphere OpenMC paired-comparison protocol and uses it to quantify local
homogenisation-bias ranges for the baseline composition and a 10--60 vol.\% Be
multiplier sweep.  Second, it performs a PF sweep and independent PF=0.618
hard-sphere seed replicates to separate topology sensitivity and
packing-to-packing variability from Monte Carlo tally uncertainty.  Third, it
uses 50-group spectra and mesh tallies to connect the global tally shifts to
fast-neutron streaming, threshold Be multiplication and spatial energy-deposition
redistribution.  The resulting assessment route identifies when homogenisation
is likely adequate and when explicit geometry or a local correction is required.

The remainder of the paper is organised as follows.  Section~2 describes the
homogenised and geometry-resolved OpenMC models, strict hard-sphere RVE
generation, tally definitions and simulation settings.  Section~3 presents the
microstructural characterisation, baseline homogenised--resolved comparison,
Be-fraction mechanism sweep, packing-fraction sensitivity, independent-seed
repeatability and spectral/spatial diagnostics.  Section~4 discusses the
relation to previous pebble-bed neutronics studies, the role of boundary
conditions and the engineering interpretation of the RVE results.  Section~5
summarises the main conclusions and identifies the validation steps required
before transferring local RVE corrections to full blanket calculations.

% ============================================================
\section{Methods}
% ============================================================

The calculation was designed as a paired comparison between two descriptions of
the same material inventory: a volume-homogenised breeder-bed material and a
geometry-resolved pebble bed.  The audited production cases use strict
non-overlapping hard-sphere geometries.  This design ensures that changes in
neutronic response are attributable to geometry, packing fraction or material
labelling rather than to unintended changes in material inventory.

\subsection{Homogenised model}

The homogenised model represents the pebble bed and inter-pebble voids as a
single uniform material.  Its nuclide density is the volume-weighted average of
the constituent solids and the He purge gas at a nominal packing fraction
(\textit{PF}) of approximately 0.62; the audited baseline OpenMC comparison uses
PF=0.618.  The density and atomic fractions follow directly from
mass conservation.  For each nuclide $j$,
%
\begin{equation}
n_j^{\rm (hom)} = {\rm PF}\cdot\Bigl[f_{\rm Li_4SiO_4}\,n_j^{\rm (b)}
+ f_{\rm Be}\,n_j^{\rm (Be)}\Bigr]
+ (1-{\rm PF})\,n_j^{\rm (He)},
\label{eq:hom}
\end{equation}
%
where $f_{\rm Li_4SiO_4}=0.60$ and $f_{\rm Be}=0.40$ are the solid volume
fractions of breeder and multiplier, respectively, and $n_j^{(\cdot)}$ are the
nuclide densities in the corresponding pure materials.


\subsection{Strict hard-sphere RVE generation}

The resolved geometries used for the audited production results were generated
as strict non-overlapping hard-sphere RVEs.  The baseline bed contains 800
spherical pebbles in a cubic RVE.  A dense parent packing was generated and then
perturbed while enforcing hard-sphere non-overlap and wall clearance.  The
resulting particle lists were audited before OpenMC import by checking the RVE
volume, total sphere volume, per-material sphere volume, packing fraction,
nearest-neighbour gaps and overlap count.  Production geometries with any
overlapping pair were excluded from the final evidence chain.

Particles were assigned as Li$_4$SiO$_4$ breeder or Be multiplier pebbles
according to the prescribed solid volume fraction.  The baseline inventory used
60 vol.\% breeder and 40 vol.\% Be, with 90 \% $^6$Li enrichment.  For the
Be-fraction mechanism sweep, the particle positions and radii were kept fixed
and only the breeder/Be labels were reassigned.  This separates the effect of
multiplier inventory from changes in packing topology.  For the PF sweep,
strict non-overlapping hard-sphere RVEs were generated at PF=0.56, 0.58, 0.60,
0.61 and 0.618.  These cases are treated as controlled topology perturbations,
not as physical thermal-expansion or irradiation-swelling calculations.

\subsection{Independent realisations and geometry validity}

To assess whether the PF=0.618 result was tied to one particle arrangement,
three independent strict hard-sphere realisations were generated and each was
transported as a paired homogenised/resolved OpenMC calculation.  The
independent-seed study quantifies packing-to-packing variability separately
from OpenMC batchwise statistical uncertainty.  All production hard-sphere
particle lists used in the PF sweep and seed-repeatability study have zero
overlap pairs in the pre-transport geometry audit.  Direct radius-scaled
thermal-expansion or swelling geometries are not used for the final physical
claims unless they can be regenerated or relaxed into strictly non-overlapping
geometries while conserving material inventory.

\subsection{OpenMC geometry and materials}

OpenMC~\cite{romano2015} was used for all neutron-transport calculations.
Each pebble in the CSV particle listing was converted into an individual
\texttt{openmc.Sphere} CSG region.  The inter-pebble void cell was the Boolean
complement of all 800 sphere interiors within the RVE bounding box, giving a
well-defined geometry of 801 cells (800~pebbles + 1~void).  Boundary conditions
were reflective in $X$ and $Y$ (infinite sheet approximation) and vacuum at
$\pm Z$.  The neutron source was an independent fixed source with strength
1.0 in OpenMC source-particle normalisation.  Its spatial distribution was a
planar box over $0\leq x\leq1.6$\,cm and $0\leq y\leq1.6$\,cm at $z=0$; the
angular distribution was monodirectional along $+z$; and the energy
distribution was discrete at 14.07\,MeV, representing D--T source neutrons.

\begin{table}[ht]
\centering
\caption{Material definitions for the OpenMC neutron-transport calculations.
The table lists the mass density, nuclide composition and temperature assigned
to the Li$_4$SiO$_4$ breeder pebbles, Be multiplier pebbles and He purge-gas
void phase in both the homogenised and geometry-resolved representative volume
element (RVE) models.  The main PF and composition sweeps used Evaluated
Nuclear Data File, release B-VII.1 (ENDF/B-VII.1), data distributed by the
National Nuclear Data Center (NNDC).  Dedicated nuclear-data sensitivity pairs
used ENDF/B-VIII.0 and FENDL-3.2c as described in Section~2.5.}
\label{tab:materials}
\small
\setlength{\tabcolsep}{3pt}
\begin{tabular}{>{\raggedright\arraybackslash}p{2.5cm}>{\raggedright\arraybackslash}p{1.9cm}>{\raggedright\arraybackslash}p{4.0cm}p{0.9cm}}
\toprule
\textbf{Material} & \textbf{Density (g/cm$^3$)} & \textbf{Nuclides} & \textbf{$T$ (K)} \\
\midrule
Li$_4$SiO$_4$ breeder & 2.40  & $^6$Li, $^7$Li (90 \% enriched), $^{28}$Si, $^{16}$O & 823 \\
Be multiplier         & 1.85  & $^9$Be  & 823 \\
He void               & 0.001 & $^4$He  & 600 \\
\bottomrule
\end{tabular}
\end{table}

\subsection{Tallies}

Global tallies were scored in fixed-source normalisation per source neutron.
The scalar tallies were \texttt{TBR\_global}, using the OpenMC score
\texttt{(n,t)} for $^6$Li; \texttt{Be\_n2n}, using \texttt{(n,2n)} for
$^9$Be; \texttt{heating\_total}, using the OpenMC \texttt{heating} score;
\texttt{flux\_total}, using the track-length \texttt{flux} score; total and
lithium-isotope-resolved \texttt{H3-production} tallies where available; and
\texttt{dpa\_total}, using \texttt{damage-energy}.  The legacy tally name
\texttt{TBR\_global} is retained in the data files, but the reported quantity
is the local RVE $^6$Li(n,t) tritium-production response per source neutron
(TPR$_6$), not a full-blanket design TBR.  Total H3-production response is
reported separately and is also a local RVE response per source neutron rather
than a full-blanket TBR.  OpenMC \texttt{heating} and \texttt{damage-energy}
scores are native eV/source
quantities; heating values are converted to MeV/source for reporting by
multiplying by $10^{-6}$, whereas damage-energy is retained in eV/source.
The unfiltered flux tally is reported as the OpenMC track-length flux tally
(particle-cm/source) and is not volume-normalised.  When photon transport is
enabled, an unfiltered OpenMC flux score includes all transported particle
types; therefore, the photon-transport sensitivity run is used for heating-score
interpretation, whereas its unfiltered flux value is not interpreted as neutron
flux.  Spatial tallies used an
$8\times8\times8$ regular Cartesian mesh over the RVE with the same
\texttt{(n,t)}, \texttt{heating}, \texttt{flux}, \texttt{(n,2n)} and
\texttt{damage-energy} scores.  Mesh heating is converted to MeV/source per
mesh bin when reported; mesh flux is likewise a track-length tally per source
particle unless explicitly divided by mesh-cell volume.

\subsection{Simulation parameters and convergence}

The archived and newly generated statepoints were confirmed to use OpenMC
fixed-source mode rather than eigenvalue mode.  The audited hard-sphere PF
sweep used 80 fixed-source batches with $3\times10^4$ source particles per
batch.  The independent-seed repeatability study used 60 fixed-source batches
with $2\times10^4$ source particles per batch.  The spectrum diagnostic used 80
fixed-source batches with $3\times10^4$ source particles per batch.  The
all-reflective boundary sensitivity pair used 10 fixed-source batches with
$5\times10^3$ source particles per batch because suppressing axial leakage
substantially lengthened particle histories in the explicit RVE.  The
photon-transport heating-sensitivity pair used 80 fixed-source batches with
$3\times10^4$ source particles per batch.  Tally means and
one-standard-deviation uncertainties are therefore batch statistics over
fixed-source realisations.
Some archived XML files contain legacy \texttt{inactive} values of 50 for the
main RVE cases and 20 for the Be-fraction scan; these entries are not
eigenvalue inactive cycles and are not associated with fission-source
convergence.  Photon transport was disabled in the main neutron-mode PF sweep
and enabled only in the dedicated heating-sensitivity pair.  The random seed was
42, and the statepoints were generated with OpenMC 0.15.0.  The main RVE cases
used ENDF/B-VII.1 NNDC HDF5 data.  Two dedicated nuclear-data sensitivity
pairs were rerun for the PF=0.618 axial-vacuum case.  The first used a minimal
ENDF/B-VIII.0 neutron HDF5 library~\cite{brown2018endfb8} generated from
official NNDC ENDF/B-VIII.0 source files using NJOY2016.68~\cite{njoy2016}
and OpenMC 0.15.0.  This minimal library
includes the nuclides required by the present model ($^6$Li, $^7$Li,
$^9$Be, $^{16}$O, $^{28}$Si and $^4$He) at 293.6, 600 and 900 K, but it does
not include photon data and is therefore used only for the neutron-mode
axial-vacuum sensitivity pair.  The second used a FENDL-3.2c neutron ACE
library~\cite{iaeaFendl32c,schnabel2024fendl} downloaded locally from the IAEA, transferred to
the remote workstation and converted to OpenMC HDF5 with
\texttt{openmc\_data} 2.4.0~\cite{openmcData} and OpenMC 0.15.0.  The newly
generated boundary, photon-transport, ENDF/B-VIII.0 and
FENDL-3.2c runs retained captured OpenMC standard output; no lost-particle
warning was found.  The retained legacy PF-sweep output files contain tally
summaries but not complete OpenMC standard output.  The
resulting relative statistical uncertainty of the RVE-level $^6$Li(n,t) tally
was 0.21--0.24 \% across the main PF sweep.

\subsection{Benchmark context and code verification}

The present RVE is not intended to reproduce the full geometry of a specific
integral benchmark.  Instead, the calculation is anchored to benchmarked
components of fusion blanket neutronics and then uses an internally paired
comparison to isolate pebble-scale topology.  The FNG-HCPB experiment provides
the closest experimental analogue because it combines 14 MeV D--T source
irradiation, Be multiplier material, Li-bearing breeder regions, tritium
production measurements and neutron-flux diagnostics in a helium-cooled pebble
bed test blanket module mock-up~\cite{batistoni2007}.  The WCCB D--T
irradiation experiment of Zhu et al.~\cite{zhu2021} provides an additional
ceramic breeder benchmark for tritium-production prediction in a
CFETR-relevant blanket mock-up.  At the code level, Valentine et al.~\cite{valentine2022}
converted several fusion benchmark problems, including FNG-HCPB responses, to
OpenMC and Serpent and compared them with established Monte Carlo workflows and
available experimental data.  Those results support the use of OpenMC for the
source-normalised reaction-rate quantities considered here.

In addition, a local executable benchmark was run for this study using a
constructive-solid-geometry reproduction of the OKTAVIAN aluminium-sphere
experiment.  The model consisted of a 20 cm aluminium sphere, a 0.2 cm
stainless-steel shell, a central D--T source spectrum and a 134-group neutron
surface-current tally at the outer boundary.  The calculation used OpenMC
0.15.0, ENDF/B-VII.1 data, 50 fixed-source batches and $2\times10^5$ source
particles per batch.  After normalising the surface current to
particle/source/cm$^2$, the 134-group leakage spectrum gave a median C/E of
1.213, a geometric-mean C/E of 1.046, 85.7 \% of valid groups within a factor
of two of experiment and a log-spectrum correlation coefficient of 0.823.  The
broad-group C/E values were 1.413 for 0.10--1 MeV, 1.184 for 1--5 MeV, 1.194
for 5--10 MeV, 0.724 for 10--15 MeV and 0.924 for 15--20.66 MeV.  This result
shows that the OpenMC installation and tally normalisation used in the present
work reproduce the overall magnitude and energy dependence of an established
14 MeV fusion-neutronics leakage benchmark.

As a closer blanket-specific benchmark target, a reconstruction of
the JAEA/FNS Li$_2$O pebble-bed mock-up of Sato et al.~\cite{sato2006} was
also implemented with the same OpenMC post-processing workflow.  The model used
the published Be/F82H/Li$_2$O slab dimensions, natural-Li Li$_2$O packing
fraction of 57.8 \%, two aluminium detector channels and eight detector
segments per channel.  Detector H3-production tallies were normalised by the
number of $^6$Li atoms in each detector segment and compared with digitised
TPR curves from the published figures.  With an isotropic monoenergetic
D--T source, the model reproduced the axial U-shaped tritium-production
profile, with shape correlations of 0.897 and 0.943 for the two detector
channels, but underpredicted the absolute TPR magnitude by about one third
(mean OpenMC/experiment = 0.666 and 0.670).  A source-angle sensitivity
calculation then bracketed the experiment: a forward-hemisphere source
overpredicted the mean TPR by 34--37 \%, whereas a mixed backward/forward
hemisphere source with a forward-probability parameter of 0.741 gave mean
calculated-to-experimental ratios of 0.968 and 1.014 for detector channels 1
and 2, respectively.  The corresponding channel-wise shape correlations were
0.969 and 0.904.  These results show that the present OpenMC geometry,
detector tally and $^6$Li inventory-normalisation chain can reproduce the
order of magnitude and axial shape of an independent Li$_2$O pebble-bed TPR
experiment when the source angular bias is represented.  However, this remains
a source-angle-calibrated screening comparison rather than a formal C/E
validation, because the published paper does not provide the full FNS
energy--angle source distribution and the present reconstruction still uses a
homogenised, rather than explicitly repeated, Li$_2$O pebble-bed layer.

\begin{figure}[H]
\centering
\includegraphics[width=\textwidth]{oktavian_al_benchmark_comparison.pdf}
\caption{Comparison of the present OKTAVIAN-Al OpenMC benchmark calculation
with experimental leakage data and public reference Monte Carlo results.
Panel (a) compares broad-group calculated-to-experimental (C/E) ratios for the
present OpenMC 0.15.0 calculation, a reference OpenMC 0.14.0 calculation and a
reference MCNP 6.1.1b calculation.  The grey band denotes agreement within a
factor of two.  Panel (b) compares integral agreement metrics for the same
datasets.  The reference OpenMC and MCNP data are taken from the public
OKTAVIAN benchmark database used by the OpenMC Fusion Benchmarks project.}
\label{fig:oktavian_benchmark}
\end{figure}

\begin{figure}[H]
\centering
\includegraphics[width=\textwidth]{li2o_fns_mixed_p0741_openmc_vs_sato_digitized.pdf}
\caption{Source-angle-calibrated OpenMC reconstruction of the JAEA/FNS
Li$_2$O pebble-bed mock-up compared with digitised tritium-production-rate
(TPR) data from Sato et al.~\cite{sato2006}.  The OpenMC model uses the
published Be/F82H/Li$_2$O slab dimensions, natural-Li Li$_2$O packing fraction
of 57.8 \%, two aluminium detector channels and eight detector segments per
channel.  Panel (a) compares detector-cell H3-production tallies normalised per
$^6$Li atom per source neutron with the digitised experimental curves.  Panel
(b) shows the OpenMC-to-experiment ratio for each detector segment; the grey
band marks the 0.95--1.15 interval reported by Sato et al. for the original
homogeneous-model C/E comparison, except at the rear boundary.  Error bars
denote one standard deviation (1$\sigma$) from the OpenMC detector tallies.
The source angular distribution is a calibrated mixed backward/forward
hemisphere approximation and is used here only as a screening benchmark, not as
a replacement for the complete FNS source description.}
\label{fig:li2o_fns_benchmark}
\end{figure}

Beyond this code-verification benchmark, the present RVE workflow was checked
by four RVE-specific consistency tests.  First, the homogenised and resolved
cases use the
same material inventory, so the reported differences are not caused by changes
in Li$_4$SiO$_4$, Be or He atom densities.  Second, all hard-sphere production
geometries were audited for packing fraction, wall clearance and zero sphere
overlap before OpenMC import.  Third, tally scores, units and uncertainty
propagation were audited directly from the available statepoint or raw result
files.  Fourth, independent PF=0.618 hard-sphere realisations reproduced the
same sign and order of the resolved--homogenised shifts.  Thus the external
benchmarks support the transport-code and blanket-physics basis, and the
Li$_2$O/FNS screening comparison checks the tritium-production tally
normalisation against an independent pebble-bed experiment.  The paired RVE
audits then support the interpretation that the present conclusions arise from
pebble-scale void topology rather than from inventory or numerical bookkeeping
artefacts.  A formal validation-grade extension would require the complete FNS
energy--angle source description or an independent Li$_4$SiO$_4$/Be benchmark
with published geometry, source and tally definitions.

\subsection{Case matrix}

Table~\ref{tab:cases} lists the calculation groups used in the audited
hard-sphere analysis.  The previously explored direct radius-scaling
deformation cases are not used as production evidence because such geometries
must first pass strict non-overlap and material-inventory-conservation audits.

\begin{table}[ht]
\centering
\caption{Audited simulation groups for the paired homogenised and
geometry-resolved RVE calculations.  PF denotes packing fraction.  The table
lists the main ENDF/B-VII.1 axial-vacuum production and diagnostic groups; the
boundary-condition, photon-transport and nuclear-data sensitivity pairs are
reported separately in Section~3.5.  Each paired calculation uses the same
material inventory and 14.07 MeV deuterium--tritium source.  The PF sweep and
independent-seed calculations use strict non-overlapping hard-sphere
geometries.}
\label{tab:cases}
\small
\setlength{\tabcolsep}{3pt}
\begin{tabular}{>{\raggedright\arraybackslash}p{2.6cm}>{\raggedright\arraybackslash}p{7.2cm}>{\raggedright\arraybackslash}p{1.2cm}}
\toprule
\textbf{Group} & \textbf{Purpose} & \textbf{PF or range} \\
\midrule
Baseline pair & Homogenised versus geometry-resolved hard-sphere RVE & 0.618 \\
PF sweep & Topology sensitivity in paired hard-sphere RVEs & 0.560--0.618 \\
Be-fraction sweep & Multiplier-inventory sensitivity at fixed topology & $\approx0.62$ \\
Seed replicates & Independent hard-sphere realisations for repeatability & 0.618 \\
Spectrum diagnostic & 50-group flux comparison for mechanism analysis & 0.618 \\
Mesh diagnostic & Mid-plane spatial redistribution analysis & 0.618 \\
\bottomrule
\end{tabular}
\end{table}

% ============================================================
\section{Results}
% ============================================================

\subsection{The generated RVE reproduces a dense random pebble bed}

The generated packing provides a reproducible microstructural reference for the
neutronic comparisons.  The nominal RVE reaches PF\,$=0.6200$, whereas the
audited baseline transport comparison uses PF=0.618 after strict hard-sphere
geometry filtering and material-inventory matching.  The interior axial
packing-fraction profile that is nearly uniform and boundary perturbations
confined to approximately one pebble diameter at the $\pm Z$ faces
(Fig.~\ref{fig:microstructure}).  The coordination-number (CN) distribution peaks at
CN$=7$ with a mean of 6.6, and the specific surface area ($S_v$) is
$S_v=20.4$\,cm$^{-1}$.  These quantities are consistent with dense random
packings of polydisperse ceramic pebbles and support the use of this RVE as a
controlled local geometry rather than as an arbitrary sphere cloud.

The combination of PF, radius distribution, CN distribution and specific
surface area also indicates that the selected packing is a typical dense
configuration for a polydisperse ceramic pebble bed within the present parameter
space.  Independent random realisations with the same nominal PF, particle-size
distribution and Li$_4$SiO$_4$/Be volume fractions would alter the exact
locations of contacts, void throats and local first-collision sites.  However,
they would be expected to preserve the same statistical void topology, mean
coordination environment and interfacial area scale.  Such sample-to-sample
differences may therefore introduce moderate stochastic fluctuations in local
mesh tallies, but they are unlikely to reverse the overall trend or order of
magnitude of the homogenisation bias, which is controlled by the presence of an
explicit connected He void topology rather than by a particular particle
arrangement.

The near-wall packing perturbation should be distinguished from the bulk-bed
region sampled by the present homogenisation comparison.  The disturbed layer is
limited to roughly one pebble diameter near the axial boundaries, whereas most
of the RVE retains a nearly constant PF close to 0.62.  Its influence on the
volume-integrated OpenMC tallies is therefore limited compared with the
dominant contribution from the interior pebble network.  Accordingly, this RVE
is used to assess homogenisation bias in the approximately uniform bulk region
of a pebble bed, not wall-boundary-layer behaviour.

The use of a finite 800-particle RVE is a limitation of the present study.  It
provides a controlled microstructural realisation for isolating the effect of
explicit Li$_4$SiO$_4$/Be/He topology.  To reduce the risk that the result is tied
to one particle arrangement, three independent PF=0.618 strict hard-sphere
realisations were also transported and compared with their matched homogenised
models.  Future work should extend this repeatability check to larger RVEs and,
where mechanically realistic packed beds are required, to independently
validated DEM or experimental particle configurations.

\begin{figure}[H]
\centering
\includegraphics[width=\textwidth]{fig1_microstructure.pdf}
\caption{Strict non-overlapping hard-sphere pebble-bed microstructure used
to construct the geometry-resolved RVE.  The data are taken from the audited
particle list exported to the OpenMC constructive solid geometry model.  Panel
(a) shows the $x$--$z$ projection of Li$_4$SiO$_4$ breeder and Be multiplier
pebbles; both spatial axes are in cm.  Panel (b) shows the local axial packing
fraction (PF) as a function of $z$ (cm), with the dashed line indicating the
nominal mean PF of 0.620.  The audited baseline OpenMC comparison uses
PF=0.618.  Panel (c) shows the pebble-radius distribution in mm, and
panel (d) shows the dimensionless coordination number (CN), defined as the
number of mechanical contacts per pebble.  These are deterministic geometry
metrics; no Monte Carlo statistical uncertainty is associated with this
figure.}
\label{fig:microstructure}
\end{figure}

\subsection{Resolved geometry changes the global reaction balance}

Replacing the homogenised bed with explicit pebbles produces a statistically
significant change in all principal tallies.  Relative to the
material-equivalent homogenised model at PF=0.618, the resolved baseline
increases TPR$_6$ by $+17.47\pm0.71$ \%, total H3-production response by
$+29.13\pm0.25$ \% and $^9$Be(n,2n) by $+45.80\pm0.11$ \%, while decreasing the
local heating tally by $-18.49\pm0.04$ \% and increasing the track-length flux by
$+23.08\pm0.09$ \%.  These differences are substantially larger than the
propagated Monte Carlo statistical uncertainties.  The corresponding absolute
homogenised and resolved tally values are retained in the supporting data rather
than repeated as a main-text table, because the same relative shifts are shown
graphically in the PF-sensitivity analysis.

The comparison is controlled by construction.  All 13 material-inventory
quantities, including $^6$Li, $^7$Li, $^9$Be, $^{28}$Si, $^{16}$O and $^4$He
atom counts, solid and void volumes, and total masses, agree to $<0.02$ \%
between the two models.  The tally shifts therefore reflect the spatial
arrangement of the same three-phase Li$_4$SiO$_4$/Be/He medium rather than a
change in elemental inventory.  In the homogenised model, the He void fraction
is averaged into a uniform mixture.  A 14.07\,MeV D--T neutron therefore samples
breeder, multiplier and purge gas through an effectively continuous macroscopic
cross section.  In the resolved model, the inter-pebble He phase forms connected
low-collision-probability channels.  These channels increase the probability
that a source neutron travels several pebble diameters before its first
collision, shifting the first-collision distribution toward discrete Be and
breeder interfaces.

This altered collision ordering explains why the relative increase in
$^9$Be(n,2n) is much larger than the increase in TPR$_6$.  The $^9$Be(n,2n)
reaction is driven by the fast part of the spectrum and is therefore highly
sensitive to whether 14-MeV neutrons retain sufficient energy when they enter
Be-rich regions.  Explicit void channels preserve high-energy flight paths and
increase the fast-neutron path length available for threshold multiplication in
Be.  TPR$_6$, by contrast, is governed by the integrated production of tritium in
$^6$Li after neutron multiplication, scattering and partial moderation.  The
$^6$Li(n,t)$^4$He reaction samples a broader transport history than the
threshold multiplication reaction.  Additional secondary neutrons generated in
Be must still be transported, moderated and absorbed in breeder material before
contributing to TPR$_6$.  Consequently, the Be multiplication rate responds more
strongly to explicit void topology, whereas TPR$_6$ shows a smaller but still
statistically significant increase.

The reduction in the local OpenMC heating-score tally has the same transport origin.
It does not indicate a lower amount of reacting material, because the inventory
is conserved.  Instead, the resolved void topology redistributes energy
deposition in space and energy.  A larger fraction of the
source-neutron population streams through He channels, undergoes multiplication
or scattering at displaced Be/breeder locations, and contributes to leakage
through the vacuum faces of the open RVE.  Energy that would be deposited
locally in the homogenised mixture is therefore partly shifted to different
positions, reaction channels or leakage histories.  The simultaneous increase
in total flux and decrease in local heating is thus physically consistent with
longer neutron residence paths and spatially displaced energy deposition, not
with an inconsistency in the heating tally.

\subsection{The geometry effect persists across multiplier inventory}

The baseline homogenisation bias remains positive over a broad range of
breeder-to-multiplier ratios.  To test composition sensitivity, the explicit
particle coordinates and radii were kept fixed while the solid-pebble inventory
was relabelled from 10 to 60\,vol.\% Be.  Each relabelled resolved case was
paired with a homogenised case having the same total solid inventory and the
same axial vacuum boundary condition.  This protocol isolates the effect of
void topology from the effect of changing the total amount of Be.

Across the full sweep, the resolved model enhances TPR$_6$ by
$+13.62$--$+17.91$ \% (Fig.~\ref{fig:mechanism_sweep}).  The $^9$Be(n,2n)
enhancement is larger but decreases monotonically from $+52.30\pm0.11$ \% to
$+40.05\pm0.07$ \% as the Be inventory increases, indicating partial saturation
of the multiplier response.  By contrast, the
total flux enhancement remains positive and slowly increases from
$+21.84\pm0.06$ \% to $+23.70\pm0.07$ \%, whereas the OpenMC heating-score response remains
lower in the resolved geometry ($-21.20\pm0.03$ \% to $-17.33\pm0.03$ \%).
Thus, the resolved-geometry
effect is not a single-composition anomaly, but a systematic transport effect
whose partitioning among breeding, multiplication and heating depends on
multiplier inventory.

The monotonic decrease in the relative $^9$Be(n,2n) enhancement reflects the
evolution of the fast-neutron field as the multiplier fraction increases.  The
$^9$Be(n,2n) channel is a threshold reaction controlled by the portion of the
spectrum that remains above the reaction threshold after 14.07\,MeV source
neutrons stream through He voids and scatter in the solid pebbles.  At low Be
fraction, the homogenised model dilutes Be uniformly through the RVE.  The
resolved model instead allows high-energy neutrons to stream through connected
He pores and intersect discrete Be pebbles before substantial energy
degradation.  Explicit geometry therefore produces a large relative increase in
the probability that a fast neutron samples a Be target while it still has
sufficient energy for multiplication.  As the Be fraction increases, both the
homogenised and resolved models already provide a high probability of Be
encounter along a typical neutron trajectory.  The additional benefit of
explicit void topology then becomes smaller in relative terms, even though the
absolute multiplication probability increases with Be inventory.

This behaviour indicates multiplier-reaction saturation rather than a loss of
the geometric effect.  Increasing the number of Be pebbles raises the material
concentration available for $^9$Be(n,2n), but it also increases competition
among neighbouring Be regions for the same population of uncollided or weakly
moderated fast neutrons.  Once streamed neutrons are already likely to collide
with Be within their first few collisions, additional Be mainly replaces
breeder or He space along trajectories that have partly lost energy.  These
later collisions contribute less efficiently to the threshold reaction because
the spectrum has been softened by elastic and inelastic scattering.  Resolved
pebble contacts and finite Be domains also introduce a self-shielding-like
effect.  The outer portions of Be pebbles preferentially intercept the
fast-neutron current, whereas inner or downstream Be volume is sampled after
attenuation and energy degradation.  Thus, the material-concentration effect
increases the available number of Be target atoms, while the geometry effect
controls how much high-energy path length is exposed to those targets.  The
decreasing relative gain from $+52.30$ \% to $+40.05$ \% indicates that the latter
effect approaches a saturated regime as the accessible fast-neutron current
becomes the limiting quantity.

The coupled trends in TPR$_6$ and flux are consistent with this interpretation.
The total flux enhancement remains nearly constant over the Be-fraction sweep,
showing that the principal topological change is the redistribution of neutron
free paths by the connected He void topology rather than a strong change in the
overall neutron population.  The reaction-specific responses then depend on how
this redistributed flux samples energy-dependent cross sections.  The
$^9$Be(n,2n) response decreases in relative magnitude because the threshold
multiplier channel is limited by the availability of fast neutrons in Be-rich
regions.  TPR$_6$ varies more weakly because tritium production integrates
multiplication, scattering, moderation and absorption in $^6$Li-bearing
Li$_4$SiO$_4$ pebbles.  Secondary neutrons generated in Be can increase the
neutron supply to the breeder, but the concurrent reduction in breeder fraction
and the requirement for transport into Li$_4$SiO$_4$ moderate the TPR$_6$ response.
The Be-fraction sweep therefore separates a material-inventory effect, which
changes the number of multiplier targets, from a persistent topology effect,
which controls neutron streaming, first-collision ordering and the exposure of
Be to the fast part of the D--T spectrum.

\begin{figure}[H]
\centering
\includegraphics[width=\textwidth]{mechanism_pf62_be_fraction_shifts.pdf}
\caption{Be multiplier-fraction sweep used to evaluate the composition
dependence of the homogenisation bias.  Each point compares a geometry-resolved
RVE with a matched homogenised RVE having the same total material inventory,
the same 14.07 MeV source and the same axial vacuum boundary condition.  Panel
(a) shows the resolved--homogenised relative shift, in percent, for TPR$_6$,
$^9$Be(n,2n), OpenMC heating-score response and total flux as functions of the Be
fraction in the solid pebble inventory (vol.\%).  Panel (b) gives the change in
each relative shift between the 10 \% and 60 \% Be cases, expressed in
percentage points (pp), to quantify the sweep amplitude.  Error bars denote one
standard deviation (1$\sigma$) propagated from the paired OpenMC tallies; all
tallies use ENDF/B-VII.1 cross-section data.}
\label{fig:mechanism_sweep}
\end{figure}


\subsection{Packing-fraction sensitivity in strict hard-sphere RVEs}

Packing-fraction sensitivity was evaluated with strict non-overlapping
hard-sphere RVEs at PF values of 0.56, 0.58, 0.60, 0.61 and 0.618.  This
sweep is interpreted as a controlled topology perturbation rather than as a
thermal-expansion or irradiation-swelling model.  For each PF, the homogenised
and resolved calculations use the same material inventory and source/boundary
conditions, so the resolved--homogenised difference remains a topology-induced
bias.

The PF sweep shows a systematic attenuation of the resolved-geometry
enhancement as the packing becomes denser (Fig.~\ref{fig:hardsphere_pf_sweep}).
The TPR$_6$ shift decreases from $+22.92\pm0.85$ \% at PF=0.56 to
$+17.47\pm0.71$ \% at PF=0.618.  The total H3-production response shift
decreases from $+34.88\pm0.29$ \% to $+29.13\pm0.25$ \%, and the
$^9$Be(n,2n) shift decreases from $+51.67\pm0.11$ \% to
$+45.80\pm0.11$ \%.  The heating shift weakens from $-21.99\pm0.04$ \% to
$-18.49\pm0.04$ \%, whereas the flux shift remains positive and nearly
constant, changing from $+22.39\pm0.09$ \% to $+23.08\pm0.09$ \%.  These
trends indicate that the explicit-geometry effect is strongest when connected
He void paths occupy a larger fraction of the RVE, because fast neutrons can
stream farther before sampling solid material.  As PF increases, the void
network narrows and the first-collision distribution of the resolved model
moves closer to the volume-averaged collision probability of the homogenised
model.  The homogenisation bias therefore decreases but does not vanish at
dense packing.

This PF dependence supports a topology-saturation interpretation for multiplier
reactions.  Increasing PF raises the solid collision probability and reduces
the relative importance of long He streaming paths.  However, the resolved
$^9$Be(n,2n) enhancement remains larger than the TPR$_6$ enhancement throughout
the sweep because the multiplier reaction remains more sensitive to the local
fast-neutron exposure of Be surfaces.  The flux shift is comparatively
insensitive to PF because the track-length tally records path-length
redistribution over the whole RVE, whereas TPR$_6$, H3 production,
$^9$Be(n,2n) and heating depend on where this redistributed flux samples
energy-dependent reaction channels.

\begin{figure}[H]
\centering
\includegraphics[width=\textwidth]{hardsphere_pf_sweep_shifts.pdf}
\caption{Packing-fraction dependence of resolved--homogenised tally shifts in
strict non-overlapping hard-sphere RVEs.  Each point represents a paired
homogenised/resolved OpenMC calculation with the same material inventory,
14.07 MeV source, reflective lateral boundaries, axial vacuum boundaries and
ENDF/B-VII.1 neutron data.  Panel (a) shows the relative differences
$(B/A-1)\times100$ \% for TPR$_6$, total H3-production response,
$^9$Be(n,2n), heating and track-length flux.  Panel (b) gives the change in
each relative shift between PF=0.560 and PF=0.618, expressed in percentage
points (pp), to highlight the densification response.  Error bars denote
propagated one-standard-deviation (1$\sigma$) tally uncertainties.}
\label{fig:hardsphere_pf_sweep}
\end{figure}

\subsection{Boundary-condition, heating-score and nuclear-data sensitivity}

The PF=0.618 baseline was extended with three targeted sensitivity checks.
First, the axial vacuum faces were
replaced by reflective boundaries while keeping the same hard-sphere particle
list, material inventory and ENDF/B-VII.1 data.  In this all-reflective local
RVE, the resolved--homogenised TPR$_6$ shift decreases from
$+17.47\pm0.71$ \% to $+15.57\pm0.72$ \%, and the total H3-production response
shift decreases from $+29.13\pm0.25$ \% to $+15.91\pm0.71$ \%.  The reduction
confirms that axial vacuum leakage amplifies part of the local breeding
response.  However, the $^9$Be(n,2n) shift remains large
($+61.52\pm0.87$ \%), showing that the threshold multiplier response is still
controlled by the exposure of discrete Be surfaces to fast neutrons even when
leakage is suppressed.  The all-reflective result should be interpreted as a
local sensitivity calculation rather than a blanket-scale reflected-boundary
design value, because the closed RVE also increases particle residence time and
does not include structural attenuation or multilayer blanket leakage paths.

Second, the PF=0.618 axial-vacuum pair was rerun with photon transport enabled.
The coupled photon-transport run gives a
heating shift of $-17.70\pm0.17$ \%, close to the neutron-mode heating shift of
$-18.49\pm0.04$ \%.  This agreement indicates that the sign and approximate
magnitude of the local heating redistribution are not artifacts of disabling
photon transport in the main neutron-mode calculations.  The reaction-rate
shifts in the photon run also remain consistent with the neutron-mode baseline,
with TPR$_6$, H3 production and $^9$Be(n,2n) shifts of
$+17.16\pm0.70$ \%, $+29.02\pm0.25$ \% and $+45.75\pm0.11$ \%, respectively.
The unfiltered flux tally from the photon-transport run is retained in the
source CSV for reproducibility but is not used as a neutron-flux metric because
OpenMC includes all transported particle types in an unfiltered flux score when
photon transport is enabled.

Third, the PF=0.618 axial-vacuum pair was repeated with generated
ENDF/B-VIII.0 and FENDL-3.2c neutron libraries.  The resolved--homogenised
shifts remain consistent with the ENDF/B-VII.1 baseline.  With ENDF/B-VIII.0,
TPR$_6$ changes from $+17.47\pm0.71$ \% to $+18.00\pm0.74$ \%, total H3
production from $+29.13\pm0.25$ \% to $+29.32\pm0.26$ \%, $^9$Be(n,2n) from
$+45.80\pm0.11$ \% to $+46.11\pm0.11$ \%, heating from $-18.49\pm0.04$ \% to
$-19.01\pm0.04$ \% and track-length flux from $+23.08\pm0.09$ \% to
$+23.14\pm0.09$ \%.  With FENDL-3.2c, the corresponding shifts are
$+17.73\pm0.73$ \%, $+29.22\pm0.26$ \%, $+45.78\pm0.11$ \%,
$-18.64\pm0.04$ \% and $+23.09\pm0.09$ \%.  The sign and magnitude of the
topology-induced local bias are therefore not specific to the ENDF/B-VII.1
evaluation for the nuclides and neutron-mode responses considered here.  These
results should be interpreted as targeted neutron-data checks, not as a
complete nuclear-data uncertainty analysis.

\begin{table}[ht]
\centering
\caption{PF=0.618 sensitivity summary for paired homogenised and geometry-resolved RVE calculations.  All values are resolved--homogenised relative shifts, $(B/A-1)\times100$ \%, with propagated one-standard-deviation (1$\sigma$) statistical uncertainty.  TPR$_6$ is the local $^6$Li(n,t) tritium-production response per source neutron.  The axial-vacuum ENDF/B-VII.1, ENDF/B-VIII.0, FENDL-3.2c and photon-transport rows used 80 fixed-source batches with 30000 particles per batch.  The all-reflective case used 10 fixed-source batches with 5000 particles per batch because the absence of axial leakage substantially lengthened resolved-geometry particle histories.  The photon-transport row is used for heating-score interpretation; its unfiltered flux tally contains both neutron and photon track-length contributions and is not interpreted as neutron flux.}
\label{tab:pf618_sensitivity}
\scriptsize
\setlength{\tabcolsep}{2pt}
\resizebox{\textwidth}{!}{%
\begin{tabular}{>{\raggedright\arraybackslash}p{3.0cm}ccccc}
\toprule
\textbf{Case} & \textbf{TPR$_6$} & \textbf{H3 production} & \textbf{$^9$Be(n,2n)} & \textbf{Heating} & \textbf{Flux} \\
\midrule
Axial vacuum, neutron mode & $+17.47\pm0.71$ & $+29.13\pm0.25$ & $+45.80\pm0.11$ & $-18.49\pm0.04$ & $+23.08\pm0.09$ \\
Axial vacuum, ENDF/B-VIII.0 & $+18.00\pm0.74$ & $+29.32\pm0.26$ & $+46.11\pm0.11$ & $-19.01\pm0.04$ & $+23.14\pm0.09$ \\
Axial vacuum, FENDL-3.2c & $+17.73\pm0.73$ & $+29.22\pm0.26$ & $+45.78\pm0.11$ & $-18.64\pm0.04$ & $+23.09\pm0.09$ \\
All-reflective, neutron mode & $+15.57\pm0.72$ & $+15.91\pm0.71$ & $+61.52\pm0.87$ & $-1.89\pm0.38$ & $+130.97\pm3.24$ \\
Axial vacuum, photon transport & $+17.16\pm0.70$ & $+29.02\pm0.25$ & $+45.75\pm0.11$ & $-17.70\pm0.17$ & $+23.30\pm0.12$$^{a}$ \\
\bottomrule
\end{tabular}
}
\begin{flushleft}
\footnotesize $^{a}$Unfiltered OpenMC flux with photon transport enabled includes neutron and photon contributions; this value is retained in the source CSV but is not used as a neutron-flux metric.
\end{flushleft}
\end{table}


\subsection{Independent hard-sphere realisations give repeatable bias trends}

The representativeness of the PF=0.618 hard-sphere result was tested with three
independent non-overlapping realisations.  The repeatability check separates
packing-to-packing variability from OpenMC tally uncertainty and directly
addresses the limitation of using a single RVE realisation.

The mean resolved--homogenised shifts over the three realisations are
$+16.78$ \% for TPR$_6$, $+28.74$ \% for total H3-production response,
$+45.69$ \% for $^9$Be(n,2n), $-18.52$ \% for heating and $+23.08$ \% for
flux (Fig.~\ref{fig:hardsphere_seed_replicates}).  The corresponding sample
standard deviations are 0.64, 0.28, 0.07, 0.04 and 0.13 percentage points,
respectively.  The sign and magnitude of the main homogenisation bias are
therefore robust to independent hard-sphere realisations within the tested
parameter set.  This does not replace a full RVE-convergence study, but it
shows that the principal mechanism is not caused by one particular particle
arrangement.

\begin{figure}[H]
\centering
\includegraphics[width=\textwidth]{hardsphere_pf618_seed_replicates.pdf}
\caption{Resolved--homogenised tally shifts for three independent PF=0.618
hard-sphere RVE realisations.  Panel (a) shows the relative shift for each
independent packing, with the horizontal ranges indicating the min--max spread
and the mean marker with sample-standard-deviation error bar.  Panel (b)
summarises the packing-to-packing sample standard deviation in percentage
points (pp).  This variability is distinct from OpenMC batchwise tally
uncertainty.}
\label{fig:hardsphere_seed_replicates}
\end{figure}


\subsection{Spectral and spatial diagnostics support the transport mechanism}

The 50-group spectrum diagnostic links the tally shifts to neutron-energy
redistribution.  In the PF=0.618 case, the resolved/homogenised flux ratio
integrated over the 0.5 eV--0.1 MeV window is 1.520, and the corresponding
ratio for 0.1--1.85 MeV is 1.549 (Fig.~\ref{fig:hardsphere_spectrum}).
Above the approximate $^9$Be(n,2n) threshold of 1.85 MeV, the integrated ratio
is 1.221.  Thus, the explicit geometry increases not only the total
track-length flux but also the population of neutrons in the fast and
intermediate energy ranges that feed multiplier reactions and subsequent
breeder absorption.

Qualitatively, the homogenised spectrum and the resolved spectrum represent two
different moderation pathways for the same 14.07 MeV D--T source.  In the
homogenised model, the source neutron continuously samples the volume-averaged
Li$_4$SiO$_4$/Be/He macroscopic cross section, so scattering, threshold
multiplication and absorption are mixed over a shortened effective transport
scale.  In the resolved model, inter-pebble He channels provide low-density
flight paths that delay the first solid collision and retain a larger
fast-neutron component near discrete Be interfaces.  This harder local
fast-neutron exposure directly supports the higher $^9$Be(n,2n) rate, while
secondary neutrons produced in Be subsequently increase the neutron supply
available for $^6$Li-bearing breeder pebbles.  The threshold-window increase
therefore provides direct mechanistic support for the observed combination of
enhanced TPR$_6$, enhanced H3-production response, enhanced multiplier reaction
rate, increased flux and reduced local heating.

\begin{figure}[H]
\centering
\includegraphics[width=\textwidth]{hardsphere_pf618_spectrum.pdf}
\caption{Fifty-group neutron spectrum diagnostic for the PF=0.618 hard-sphere
RVE.  Panel (a) shows the OpenMC track-length flux per lethargy per source
neutron for the homogenised and geometry-resolved models.  Panel (b) shows the
resolved/homogenised spectral-flux ratio for energy bins above 1 keV.  Panel
(c) summarises the integrated resolved--homogenised flux shift in selected
energy windows.  The
vertical dashed line marks the approximate $^9$Be(n,2n) threshold at 1.85 MeV.
The data are obtained from paired fixed-source OpenMC calculations with the same
material inventory, source and boundary conditions.}
\label{fig:hardsphere_spectrum}
\end{figure}

Mid-plane mesh diagnostics show that the resolved geometry redistributes local
reaction rates within the RVE.  In the PF=0.618 mid-plane, the mean mesh-bin
shifts are positive for TPR$_6$ and H3 production, negative for heating and
weakly positive for flux, but individual bins exhibit large spatial variation
because local material placement, void connectivity and tally statistics vary
at the cell scale.  These maps should therefore be interpreted as evidence of
spatial redistribution rather than as a deterministic correction field.  Their
main role is to confirm that the global shifts arise from heterogeneous
transport paths through the explicit Li$_4$SiO$_4$/Be/He topology.

\begin{figure}[H]
\centering
\includegraphics[width=\textwidth]{hardsphere_pf618_midplane_mesh_shifts.pdf}
\caption{Mid-plane mesh-bin resolved--homogenised shifts for the PF=0.618
hard-sphere RVE.  The maps are extracted from the $8\times8\times8$ Cartesian
OpenMC mesh at the RVE mid-plane.  Panels show the cell-wise relative shifts for
(a) TPR$_6$, (b) total H3-production response, (c) heating and (d) track-length
flux; panel (e) gives the corresponding distribution of mid-plane mesh-bin
shifts, with boxes denoting the interquartile ranges and vertical ticks denoting
the medians.  Spatial axes are $x$ and $y$ in cm.  Colour bars denote
$(B/A-1)\times100$ \% for each mesh bin; mesh-bin statistical uncertainties are
not plotted to preserve map readability.}
\label{fig:hardsphere_spatial}
\end{figure}

% ============================================================
\section{Discussion}
% ============================================================

The central result of this study is that homogenisation can change the local
reaction balance of a breeder-bed RVE even when material inventory is conserved
to high precision.  The effect is not limited to TPR$_6$.  $^9$Be(n,2n), total flux
and local heating change in different directions, indicating a
transport-history effect rather than a simple material-density effect.  The
physical origin is
the replacement of a connected three-phase void topology by a uniform
macroscopic cross section in the homogenised model.  This distinction matters
for coupled blanket analyses because thermal, swelling and damage calculations
often use the same Monte Carlo tallies that are affected by homogenisation.

\subsection{Relation to previous pebble-bed neutronics studies}

The positive TPR$_6$ shift found here differs from the small negative
zone-integrated Li-6 tritium-production deviation reported by Zhang and
Guo~\cite{zhang2025} for a reflected helium-cooled blanket.  In their
Li$_4$SiO$_4$ region-1 analysis, the random-packing local-bin deviations
relative to homogenisation ranged from -8.05 \% to +2.41 \%, whereas the
integrated deviation was only -0.15 \%.  This difference should not be read as
a contradiction.
Reflected-neutron effects are known to influence fusion-blanket tritium-production
interpretation~\cite{sato2008},
and the two calculations isolate different parts of the problem.  Four modelling
choices are particularly important:

\begin{enumerate}[nosep]
\item \textbf{Boundary conditions.}  Our vacuum $\pm Z$ allows streaming
      neutrons to exit, maximising the streaming effect; their reflective
      azimuthal boundaries and Be reflector layers recirculate neutrons,
      reducing net enhancement.
\item \textbf{Packing-fraction distribution.}  Our Lubachevsky--Stillinger (LS)
      algorithm gives near-constant PF
      ($\approx$0.62) throughout; their random packing exhibits $\pm$8 \%
      near-wall PF oscillations that drive compensating negative local
      deviations.
\item \textbf{Material composition.}  Our model includes 40~vol\% Be
      multiplier, whose (n,2n) reactions thermalise the spectrum and increase
      $^6$Li absorption; their Section~4 uses pure Li$_4$SiO$_4$.
\item \textbf{Source.}  Monoenergetic vs.\ Gaussian-distributed.
\end{enumerate}

The present calculation is therefore a controlled local homogenisation test,
whereas Zhang and Guo address a more blanket-like reflected configuration.
Together, the two studies indicate that explicit pebble geometry can remain
locally important even when a full-blanket global tritium-production shift is small.

\subsection{Boundary-condition scope and engineering interpretation}

The vacuum boundaries on the axial $z$ faces create a strong leakage channel
and amplify free-streaming effects through the connected He void topology.  In
this configuration, a 14.07\,MeV D--T neutron entering an inter-pebble void can
travel farther before collision, reach discrete Be or Li$_4$SiO$_4$ interfaces
with less spectral degradation, or leave the RVE through the vacuum faces.  This
boundary choice isolates the homogenisation error associated with replacing the
explicit three-phase topology by a uniform macroscopic cross section.  The
absolute homogenised--resolved differences reported here should therefore be
regarded as leakage-sensitive upper-bound estimates for this local RVE, not as
direct full-blanket design margins.

The all-reflective sensitivity calculation confirms that this boundary effect is
not merely qualitative.  Suppressing axial leakage reduces the TPR$_6$ shift from
$+17.47\pm0.71$ \% to $+15.57\pm0.72$ \% and reduces the total H3-production
shift from $+29.13\pm0.25$ \% to $+15.91\pm0.71$ \%.  The local heating shift
also weakens from $-18.49\pm0.04$ \% to $-1.89\pm0.38$ \%, indicating that
neutron return partially compensates the leakage-driven reduction in local
energy deposition.  By contrast, the $^9$Be(n,2n) shift increases to
$+61.52\pm0.87$ \%.  This behaviour is physically plausible for a closed local
RVE: reflected fast and epithermal histories increase the residence time and
allow repeated sampling of discrete Be surfaces, while the homogenised model
still distributes Be continuously through the volume.  The result therefore
shows that boundary conditions can change not only the magnitude but also the
partitioning of the homogenisation bias among breeding, multiplication, flux and
heating responses.

The engineering implication is therefore not that the numerical shifts in this
RVE can be inserted directly into a solid breeder blanket design.  Rather, the
RVE provides a controlled upper-bound test of the mechanism by which pebble-bed
topology changes neutron free path, collision ordering, spectral degradation and
reaction-site distribution.  For reflected or multilayer blanket calculations,
the appropriate transfer of this result is to retain the mechanism, not the
absolute percentage values.  A practical workflow is to use explicit
geometry-resolved RVEs as local correction or uncertainty-quantification models
inside representative blanket zones, repeat the comparison under the actual
module boundary condition or response-matrix environment, and then propagate the
resulting local homogenisation bias to full-sector Monte Carlo calculations.
Such an approach would allow the effects identified here, namely void-channel
streaming, fast-neutron exposure of Be, breeder absorption after moderation and
spatial redistribution of nuclear heating, to be evaluated under realistic
reflective neutron return, coolant-channel shielding and structural-material
attenuation.

The PF and composition sweeps are parametric local sensitivity studies.  They
show that the sign and order of the local homogenisation bias persist under
controlled hard-sphere topology and Be-inventory perturbations, but they do not
constitute a validated thermal-expansion, irradiation-swelling or full-blanket
feedback model.  The absolute RVE-level TPR$_6$, the magnitude of
leakage-driven spatial peaks and any architecture-specific breeding margin
remain outside the claims of this paper.

On this basis, a reflective or multilayer blanket environment should not be
represented by directly transferring the axial-vacuum percentages.  The new
all-reflective RVE check demonstrates sensitivity to neutron return, but it is
still a closed local cell rather than a multilayer blanket segment with
structural attenuation, coolant channels and finite leakage.  A modern
engineering application therefore requires reflected or multilayer OpenMC
calculations with the relevant blanket architecture before any quantitative
correction is applied to design.

\subsection{Engineering use of the homogenised model}

The present results can be translated into a practical screening rule for
engineering neutronics, provided that the rule is applied to local modelling
choices rather than to the final blanket breeding margin.  In regions where a
paired resolved--homogenised comparison, or an experimentally validated local
correction, indicates a TPR$_6$ deviation below about 5 \%, the homogenised model
may be acceptable for preliminary design studies, global optimisation and
parametric screening, especially when the quantity of interest is an
integral blanket response.  This threshold is not a universal safety criterion.
It is a pragmatic engineering tolerance below which the geometry-induced
uncertainty is comparable to other modelling uncertainties in early blanket
analysis.  Conversely, the present RVE gives TPR$_6$ deviations of
$+13.62$--$+17.91$ \% across the Be-fraction sweep, together with much larger
$^9$Be(n,2n) deviations and a substantial heating redistribution.  A deviation
of this order should be treated as evidence that explicit pebble geometry, or a
locally calibrated correction factor derived from resolved RVEs, is required
before the result is used in design-sensitive tritium-production or nuclear-heat
assessments.

The model choice should also depend on the structural location within the
blanket.  In the approximately uniform bulk region of a pebble bed, where the
packing fraction is close to its interior value and large void clusters are
absent, homogenisation can remain useful if the local bias is demonstrated to
fall within the acceptable range defined above.  Near-wall boundary layers,
purge-gas gaps, coolant-channel neighbourhoods, local packing defects and
regions with concentrated void connectivity require a more conservative
treatment.  In these zones, neutron streaming, first-collision ordering and
spectral hardening are governed by geometry rather than by the volume-averaged
material composition, so explicit geometry or a resolved-RVE-based correction
should be used for TPR$_6$, $^9$Be(n,2n), flux and heating tallies.

The operating state provides a second screening dimension.  For a static pebble bed
with a stable packing fraction and no significant irradiation-induced
redistribution, a homogenised model may be adequate after a one-time
resolved-RVE verification of the local bias.  If future thermal-expansion,
uniform-swelling or flux-correlated-swelling cases are used, they should first
be regenerated as non-overlapping geometries with material-inventory-conserving
density scaling.  Until such calculations are available, the present PF sweep
should be interpreted only as a controlled hard-sphere topology sensitivity
study.

Accordingly, the RVE results should be used as a correction basis and a
decision framework for full-blanket calculations, not as direct design values.
A recommended workflow is to perform full-sector calculations with homogenised
materials, identify breeder-bed zones that are close to geometry-sensitive
conditions, compute representative resolved RVEs for those zones under the
actual material composition and boundary environment, and propagate the
resulting local correction or uncertainty into the full model.  This workflow
preserves the computational efficiency of homogenised neutronics while avoiding
unqualified use of homogenisation in regions where pebble-scale topology
controls multiplication, breeding or energy deposition.

% ============================================================
\section{Conclusions and outlook}
% ============================================================

This work demonstrates that explicit pebble-scale geometry provides a
first-order local correction to homogenised neutronics in a strict hard-sphere
Li$_4$SiO$_4$/Be breeder-bed RVE.  With identical material inventory at PF=0.618,
the resolved baseline increases TPR$_6$ by $+17.47\pm0.71$ \%, total
H3-production response by $+29.13\pm0.25$ \% and $^9$Be(n,2n) by
$+45.80\pm0.11$ \%, while reducing the OpenMC heating score by
$-18.49\pm0.04$ \% and increasing track-length flux by $+23.08\pm0.09$ \%.
The hard-sphere PF sweep and Be-fraction sweep show that the bias persists
under controlled topology and multiplier-inventory perturbations, and the
independent-seed study confirms that the sign and order of the shifts are not
tied to a single particle arrangement.

The mechanism is controlled by connected He void topology.  Explicit void paths
delay solid collisions, preserve fast-neutron exposure to Be, enhance
threshold $^9$Be(n,2n) multiplication and redistribute local energy deposition.
The 50-group spectrum diagnostic supports this interpretation by showing a
resolved/homogenised flux ratio of 1.221 above the approximate $^9$Be(n,2n)
threshold.  The ENDF/B-VIII.0 and FENDL-3.2c sensitivity runs preserve the same
sign and approximate magnitude of the axial-vacuum local bias.  The
ENDF/B-VIII.0 shifts are $+18.00\pm0.74$ \% for TPR$_6$,
$+29.32\pm0.26$ \% for total H3 production, $+46.11\pm0.11$ \% for
$^9$Be(n,2n), $-19.01\pm0.04$ \% for heating and $+23.14\pm0.09$ \% for
flux.  The corresponding FENDL-3.2c shifts are $+17.73\pm0.73$ \%,
$+29.22\pm0.26$ \%, $+45.78\pm0.11$ \%, $-18.64\pm0.04$ \% and
$+23.09\pm0.09$ \%.  The all-reflective sensitivity check reduces the TPR$_6$ and
H3-production shifts to $+15.57\pm0.72$ \% and $+15.91\pm0.71$ \%, respectively,
while increasing the $^9$Be(n,2n) shift to $+61.52\pm0.87$ \%; this confirms
that boundary conditions modify the partitioning of the local bias.  A
photon-transport rerun gives a heating shift of $-17.70\pm0.17$ \%, close to the
neutron-mode value, supporting the robustness of the heating-redistribution
interpretation.  The reported quantities are source-normalised local RVE
responses, not full-blanket TBR margins.  Future work should extend the same
audited workflow to multilayer blanket environments and validated
microstructure-evolution geometries that conserve material inventory and remain
strictly non-overlapping.

\section*{Data and code availability}

The input particle lists, OpenMC input-generation scripts, post-processing
scripts, processed tally CSV files, figure source data and manuscript build
files are available at
\url{https://github.com/wangjianfttt/fusion-pebble-bed-rve-openmc-data} and
archived on Zenodo at \url{https://doi.org/10.5281/zenodo.20636518}.
Continuous-energy nuclear-data libraries are not redistributed because of their
size and licensing conditions; the exact library names, source URLs, processing
tools, checksums and local \texttt{cross\_sections.xml} paths used in the
calculations are recorded in the data package.

\section*{CRediT authorship contribution statement}

Jian Wang: Conceptualization, Methodology, Software, Investigation, Formal
analysis, Visualization, Writing -- original draft.  Haoxi Wang: Methodology,
Writing -- review and editing.  Zhiyuan Liu: Application context, Methodology,
Writing -- review and editing.  Haishun Deng: Application context, Validation
planning, Writing -- review and editing.  Wei Wen: Data curation,
Visualization, Writing -- review and editing.  Gang Shen: Supervision, Project
administration, Writing -- review and editing.

\section*{Declaration of competing interest}

The authors declare no known competing financial interests or personal
relationships that could have appeared to influence the work reported in this
paper.

\section*{Acknowledgements}

This work was supported by the Anhui Intelligent Mine Technology and Equipment
Engineering Research Center (AIMTEERC202307), the Anhui Provincial Natural
Science Foundation (2408085QA030), the Science Foundation of the Institute of
Plasma Physics, Chinese Academy of Sciences (DSJJ-2025-08), and the China
Post-doctoral Science Foundation under Grant Number 2024M753266.

% ============================================================
\begin{thebibliography}{22}

\bibitem{giancarli2018}
Giancarli L, et al.\ ITER TBM Program and associated system engineering.
\textit{Fusion Eng.\ Des.} 2018;136:815--821.
doi:10.1016/j.fusengdes.2018.04.014.

\bibitem{fischer2015}
Fischer U, et al.\ Neutronics requirements for a DEMO fusion power plant.
\textit{Fusion Eng.\ Des.} 2015;98--99:2134--2137.
doi:10.1016/j.fusengdes.2015.02.029.

\bibitem{romano2015}
Romano PK, Forget B.\ OpenMC: A state-of-the-art Monte Carlo code for research
and development. \textit{Ann.\ Nucl.\ Energy} 2015;82:90--97.
doi:10.1016/j.anucene.2014.07.048.

\bibitem{zhang2025}
Zhang S, Guo H.\ Neutronic effects of the pebble bed model on tritium breeding
calculations for fusion blankets. \textit{Fusion Eng.\ Des.} 2025;220:115377.
doi:10.1016/j.fusengdes.2025.115377.

\bibitem{sato2006}
Sato S, et al.\ Experimental study on nuclear properties of water cooled pebble
bed blanket. In: \textit{Proceedings of the 21st IAEA Fusion Energy
Conference}; 2006. Paper FT/P5-13.

\bibitem{bernal64}
Bernal JD.\ The Bakerian Lecture, 1962: The structure of liquids.
\textit{Proc.\ R.\ Soc.\ Lond.\ A} 1964;280:299--322.
doi:10.1098/rspa.1964.0147.

\bibitem{jodrey85}
Jodrey WS, Tory EM.\ Computer simulation of close random packing of equal
spheres. \textit{Phys.\ Rev.\ A} 1985;32:2347--2351.
doi:10.1103/PhysRevA.32.2347.

\bibitem{sato2008}
Sato S, et al.\ Impact of reflected neutrons on accuracy of tritium production
rate prediction in blanket mock-ups for fusion reactors.
\textit{Fusion Eng.\ Des.} 2008;83:1304--1308.
doi:10.1016/j.fusengdes.2008.08.004.

\bibitem{batistoni2007}
Batistoni P, Angelone M, Bettinali L, Carconi P, Fischer U, Kodeli I, et al.\
Neutronics experiment on a helium cooled pebble bed (HCPB) breeder blanket
mock-up. \textit{Fusion Eng.\ Des.} 2007;82:2095--2104.
doi:10.1016/j.fusengdes.2007.04.009.

\bibitem{zhu2021}
Zhu Q, Chen W, Bao J, Zhao Z, Du H, Huang K, et al.\ Experimental study on
tritium breeding in water-cooled ceramic breeder blanket mock-up under D--T
neutron irradiation conditions. \textit{Nucl.\ Fusion} 2021;61:066018.
doi:10.1088/1741-4326/abed84.

\bibitem{valentine2022}
Valentine A, Berry T, Bradnam S, Hagues J, Hodson J.\ Benchmarking of emergent
radiation transport codes for fusion neutronics applications.
\textit{Fusion Eng.\ Des.} 2022;180:113197.
doi:10.1016/j.fusengdes.2022.113197.

\bibitem{kodeli2021}
Kodeli IA, Sartori E.\ SINBAD -- Radiation shielding benchmark experiments.
\textit{Ann.\ Nucl.\ Energy} 2021;159:108254.
doi:10.1016/j.anucene.2021.108254.

\bibitem{konno2010}
Milocco A, Trkov A, Kodeli IA.\ The OKTAVIAN TOF experiments in SINBAD:
Evaluation of the experimental uncertainties.
\textit{Ann.\ Nucl.\ Energy} 2010;37:443--449.
doi:10.1016/j.anucene.2009.12.016.

\bibitem{gong2017}
Gong B, et al.\ Investigation on the packing behaviors and mechanics of
Li$_4$SiO$_4$ pebble beds by discrete element method.
\textit{Fusion Eng.\ Des.} 2017;125:551--555.
doi:10.1016/j.fusengdes.2017.04.049.

\bibitem{shen2020}
Shen Q, et al.\ Neutronics effects of homogeneous model on solid breeder
blanket of CFETR. \textit{Fusion Eng.\ Des.} 2020;160:111825.
doi:10.1016/j.fusengdes.2020.111825.

\bibitem{zhang2022}
Zhang H, et al.\ Packing characteristics of pebble beds for fusion reactors
under different friction coefficients.
\textit{Fusion Eng.\ Des.} 2022;176:113051.
doi:10.1016/j.fusengdes.2022.113051.

\bibitem{wang2020}
Wang X, et al.\ Three-dimensional validation and analyses of the optimized
CFETR HCCB blanket. \textit{Fusion Eng.\ Des.} 2020;161:111971.
doi:10.1016/j.fusengdes.2020.111971.

\bibitem{brown2018endfb8}
Brown DA, Chadwick MB, Capote R, et al.\ ENDF/B-VIII.0: The 8th major release
of the nuclear reaction data library with CIELO-project cross sections, new
standards and thermal scattering data. \textit{Nucl.\ Data Sheets}
2018;148:1--142. doi:10.1016/j.nds.2018.02.001.

\bibitem{iaeaFendl32c}
International Atomic Energy Agency Nuclear Data Section. Fusion Evaluated
Nuclear Data Library, FENDL-3.2c. 2025. Available at:
\url{https://nds.iaea.org/fendl/}; accessed 11 June 2026.

\bibitem{schnabel2024fendl}
Schnabel G, Aldama DL, Bohm T, et al.\ FENDL: A library for fusion research
and applications. \textit{Nucl.\ Data Sheets} 2024;193:1--78.
doi:10.1016/j.nds.2024.01.001.

\bibitem{njoy2016}
Macfarlane R, Muir DW, Boicourt RM, Kahler AC, Conlin JL. The NJOY Nuclear
Data Processing System, Version 2016. Los Alamos National Laboratory report
LA-UR-17-20093; 2017. doi:10.2172/1338791.

\bibitem{openmcData}
Shimwell J, et al. \texttt{openmc\_data}: tools for processing and downloading
nuclear data for OpenMC. Available at:
\url{https://github.com/openmc-data-storage/openmc_data}; accessed 11 June
2026.

\end{thebibliography}

\end{document}
